\documentclass[11pt,a4paper]{article}

\usepackage[T1]{fontenc}
\usepackage{lmodern}
\usepackage[margin=0.72in]{geometry}
\usepackage{xcolor}
\usepackage{hyperref}
\usepackage{enumitem}
\usepackage{titlesec}
\usepackage{microtype}

\definecolor{accent}{RGB}{0, 86, 179}

\hypersetup{
  colorlinks=true,
  urlcolor=black,
  linkcolor=black,
  pdfauthor={Linfei Li},
  pdftitle={Linfei Li - Curriculum Vitae}
}
\urlstyle{same}

\titleformat{\section}
  {\large\bfseries\color{accent}}
  {}{0em}{}[\color{accent}\titlerule]
\titlespacing*{\section}{0pt}{0.9em}{0.45em}

\setlength{\parindent}{0pt}
\setlength{\parskip}{0pt}
\setlist[itemize]{leftmargin=1.25em, topsep=0.25em, itemsep=0.18em, parsep=0pt, partopsep=0pt}
\setlist[enumerate]{leftmargin=1.45em, topsep=0.3em, itemsep=0.45em, parsep=0pt, partopsep=0pt}
\raggedbottom
\pagestyle{empty}

\newcommand{\twocolline}[2]{%
  \begin{minipage}[t]{0.76\textwidth}#1\end{minipage}\hfill%
  \begin{minipage}[t]{0.22\textwidth}\raggedleft #2\end{minipage}\par}
\newcommand{\datedline}[2]{\twocolline{\textbf{#1}}{#2}}
\newcommand{\roleline}[2]{\twocolline{\textit{#1}}{#2}\vspace{-0.15em}}

\begin{document}

\begin{center}
  {\LARGE\bfseries Linfei Li}\\[0.32em]
  \small
  \href{tel:+8618387400236}{(+86) 183-8740-0236}
  $\cdot$ \href{mailto:cslinfeili@tongji.edu.cn}{cslinfeili@tongji.edu.cn}
  $\cdot$ \href{https://lif314.github.io/}{lif314.github.io}
  $\cdot$ \href{https://scholar.google.com/citations?user=V-ugTYUAAAAJ&hl=en}{Google Scholar}
\end{center}
\vspace{-0.35em}

\section{EDUCATION}
\datedline{Tongji University}{Shanghai, China}
\roleline{Ph.D. Student in Computer Science, GPA: 90.67/100}{09/2023 -- 03/2029 (5.5-year program)}
Advisors: Prof. Lin Zhang and Prof. Ying Shen

\vspace{0.35em}
\datedline{Tongji University}{Shanghai, China}
\roleline{B.S. in Software Engineering, GPA: 4.79/5.0}{09/2019 -- 07/2023}
Selected coursework: Computer Vision, Computer Graphics, Data Structures and Algorithms, Programming Languages, Operating Systems, Computer Networks, Software Engineering

\section{SELECTED PUBLICATIONS}
\textit{$^*$Equal contribution.}
\begin{enumerate}[label=\arabic*.]
  \item \textbf{Linfei Li}, Ruining Hu, Lin Zhang, Zhong Wang, Fengyi Zhang, Ying Shen, Binqiang Wang, Xin Zhang. \href{https://lif314.github.io/projects/unirap/}{\textit{UniRAP: Towards Unified Part-level Physical Affordance Reasoning and Actionable Perception.}} In \textbf{NeurIPS}, 2026. (CCF-A)

  \item \textbf{Linfei Li}, Lin Zhang, Ying Shen. \href{https://lif314.github.io/projects/realvlg_r1/}{\textit{RealVLG-R1: A Large-Scale Real-World Visual-Language Grounding Benchmark for Robotic Perception and Manipulation.}} In \textbf{CVPR}, 2026. (CCF-A)

  \item Yuchuan Ding$^*$, \textbf{Linfei Li}$^*$, Lin Zhang, Ying Shen. \href{https://lif314.github.io/projects/raysup/}{\textit{RaysUp: Ultra-light Universal Feature Upsampling via Geometry-Aware Ray Representation.}} In \textbf{ECCV}, 2026. (CCF-B)

  \item \textbf{Linfei Li}, Lin Zhang, Zhong Wang, Ying Shen. \textit{SmartSplat: Feature-Smart Gaussians for Scalable Compression of Ultra-High-Resolution Images.} In \textbf{AAAI}, 2026. (CCF-A)

  \item \textbf{Linfei Li}, Lin Zhang, Zhong Wang, Fengyi Zhang, Zelin Li, Ying Shen. \href{https://lif314.github.io/projects/neaf/}{\textit{Representing Sounds as Neural Amplitude Fields: A Benchmark of Coordinate-MLPs and a Fourier Kolmogorov-Arnold Framework.}} In \textbf{AAAI}, 2025. (CCF-A)

  \item \textbf{Linfei Li}, Lin Zhang, Zhong Wang, Ying Shen. \href{https://lif314.github.io/projects/gs3lam/}{\textit{GS3LAM: Gaussian Semantic Splatting SLAM.}} In \textbf{ACM MM}, 2024. (CCF-A)

  \item Zhanbo Shi, Lin Zhang, \textbf{Linfei Li}, Ying Shen. \textit{Towards Audio-Visual Navigation in Noisy Environments: A Large-Scale Benchmark Dataset and an Architecture Considering Multiple Sound Sources.} In \textbf{AAAI}, 2025. (CCF-A)

  \item \textbf{Linfei Li}, Fengyi Zhang, Zhong Wang, Lin Zhang, Ying Shen. \textit{INR-Bench: A Unified Benchmark for Implicit Neural Representations in Multi-Domain Regression and Reconstruction.} Technical Report, 2025.
\end{enumerate}

\section{RESEARCH PROJECTS}
\datedline{National Key R\&D Program - Next-Generation AI}{Beijing, China}
\roleline{Sub-project Technical Lead: Multimodal LLM Cross-Modal Trustworthy Verification}{12/2025 -- Present}
\begin{itemize}
  \item Lead the technical roadmap and core algorithm development for trustworthy cross-modal reasoning in multimodal large language models.
\end{itemize}

\datedline{Changchun Science \& Technology Breakthrough Program}{Changchun, China}
\roleline{Sub-project Technical Lead: Multi-Sensor Fusion-Based AR Navigation}{01/2025 -- 12/2025}
\begin{itemize}
  \item Led camera-IMU-LiDAR calibration and designed the multi-sensor fusion framework for in-vehicle AR navigation.
\end{itemize}

\datedline{Tongji University Scenario Validation Program}{Shanghai, China}
\roleline{Sub-project Technical Lead: Multi-Robot Collaborative Localization and Mapping}{11/2024 -- 11/2025}
\begin{itemize}
  \item Designed the multi-robot SLAM architecture; led pose-graph optimization and dense semantic mapping from simulation through real-robot deployment.
\end{itemize}

\section{ACADEMIC SERVICE}
\textbf{Reviewer:} CVPR 2026; NeurIPS 2026; ECCV 2026; AAAI 2025--2026; ICML 2025; ACM MM 2025

\section{HONORS \& AWARDS}
\begin{itemize}
  \item Tongji University Outstanding Doctoral Student Scholarship \hfill 2025
  \item Tongji University Outstanding Doctoral Student \hfill 2025
  \item Outstanding Graduate of Tongji University \hfill 2023
  \item First Prize, National Mathematical Contest in Modeling (Shanghai Division) \hfill 2021
  \item National Endeavor Scholarship \hfill 2020, 2021, 2022
  \item Tongji Scholarship of Excellence \hfill 2020, 2021, 2022
\end{itemize}

\section{TECHNICAL SKILLS}
\textbf{Research:} Multimodal large language models, vision-language models, embodied AI and robot manipulation, visuo-tactile perception, reinforcement learning for robotics, dense prediction, 3D reconstruction\\
\textbf{Robotics:} Franka Emika robot arms, dexterous hands, tactile sensors, humanoid robots, camera-LiDAR-IMU calibration\\
\textbf{Programming:} Python, C/C++, Java \quad
\textbf{Frameworks:} PyTorch, JAX, Hugging Face, ROS/ROS2, Isaac Sim, \LaTeX

\end{document}
